\documentclass[11pt,a4paper]{article}
\usepackage[T1]{fontenc}
\usepackage[utf8]{inputenc}
\usepackage{lmodern}
\usepackage{amsmath,amssymb,amsthm,mathtools}
\usepackage[margin=28mm]{geometry}
\usepackage{microtype}
\usepackage{xcolor}
\usepackage[colorlinks=true,linkcolor=blue!45!black,citecolor=blue!45!black,urlcolor=blue!45!black]{hyperref}
\numberwithin{equation}{section}
\newtheorem{theorem}{Theorem}[section]
\newtheorem{proposition}[theorem]{Proposition}
\newtheorem{lemma}[theorem]{Lemma}
\theoremstyle{remark}
\newtheorem{remark}[theorem]{Remark}
\newcommand{\D}{\mathbb D}
\newcommand{\E}{\mathbb E}
\newcommand{\Prob}{\mathbb P}
\newcommand{\Mplus}{\mathcal M_+}
\newcommand{\Mp}{\mathcal M_{\mathrm p}}
\newcommand{\dd}{\,\mathrm d}
\newcommand{\ind}{\mathbf 1}
\setlength{\emergencystretch}{2em}
\hypersetup{pdftitle={Compactness and simplicity for the planar disk log gas},pdfauthor={}}
\title{Compactness and simplicity for the planar disk log gas}
\author{}
\date{}

\begin{document}
\maketitle
\vspace{-2em}
\begin{abstract}
For every fixed inverse temperature $\beta>0$, the unscaled counting
measures of the logarithmic gas in the open unit disk have tight laws.
Every accumulation point, even when convergence is taken in the space
of all nonnegative Radon measures, is a locally finite simple point
process. We prove a uniform Gaussian upper tail for compact counts,
derive a local two-point repulsion estimate, and show explicitly why
these bounds exclude both diffuse mass and multiple atoms in the limit.
All estimates specific to the model are proved in this note.
\end{abstract}

\section{Model, topology, and statement}

Let $\D=\{z\in\mathbb C:|z|<1\}$, let $B_r=\{z:|z|<r\}$, and
let $dA$ denote ordinary, unnormalized planar Lebesgue measure.
Fix $\beta>0$. For $N\ge1$, the labelled particles have law
\begin{equation}\label{eq:model}
 \Prob_{N,\beta}(\dd z_1\cdots\dd z_N)
 =\frac{1}{Z_{N,\beta}}|\Delta_N(z)|^\beta\prod_{i=1}^N dA(z_i),
 \qquad z\in\D^N,
 \qquad \Delta_N(z)=\prod_{1\le i<j\le N}(z_i-z_j).
\end{equation}
Here $Z_{N,\beta}$ is the integral of the numerator over $\D^N$,
and empty products equal one. The partition function is finite because
all pair distances are at most $2$, and positive because the integrand
is positive off the collision diagonals. In particular, this law is
absolutely continuous and its particles are distinct almost surely.
Write
\[
 \Xi_N=\sum_{i=1}^N\delta_{z_i},\qquad M_N(r)=\Xi_N(B_r).
\]
There is no spatial rescaling as $N$ increases.

The space $\Mplus(\D)$ consists of nonnegative Radon measures on
$\D$: their mass on every compact subset of $\D$ is finite.
It carries the \emph{vague topology}, in which $\mu_n\to\mu$ means
\[
 \mu_n(f)\longrightarrow\mu(f),\qquad
 \mu(f)=\int f\dd\mu,\qquad f\in C_c(\D).
\]
Here $C_c(\D)$ is the space of continuous functions with compact
support in $\D$. The subspace $\Mp(\D)$ consists of locally finite
counting measures, allowing integer multiplicities at an atom.
A \emph{point process} is a random element of $\Mp(\D)$; it is
\emph{simple} if the mass of each atom is one almost surely.
Convergence in law below refers to probability laws on these spaces,
not to convergence of the measures $N^{-1}\Xi_N$.

Vague convergence on the open disk permits particles to escape to the
boundary. This is essential: on the closed disk the unnormalized total
mass is $N$, so these laws could not be tight.

\begin{theorem}[Compactness and simple limits]\label{thm:main}
Fix $\beta>0$.
The laws of $(\Xi_N)_{N\ge1}$ are tight in $\Mplus(\D)$.
Every sequence of particle numbers tending to infinity therefore has
a subsequence along which $\Xi_N$ converges in law.
Every accumulation point is a locally finite simple point process.

For $0<r<1$ and every integer $m\ge2$,
\begin{equation}\label{eq:tail}
 \sup_{N\ge1}\Prob\{M_N(r)\ge m\}
 \le \min\left\{1,
 \exp\left(m-\frac{\beta}{4}m(m-1)\log\frac1r\right)\right\}.
\end{equation}
The same bound holds for $M(r)=\Xi(B_r)$ for every limiting process
$\Xi$. For every compact $K\subset\D$, $q>0$, and $t\in\mathbb R$,
\[
 \sup_N\E[\Xi_N(K)^q]<\infty,\qquad
 \sup_N\E e^{t\Xi_N(K)}<\infty,
\]
and these moments are finite for every limit.
Along any subsequence $\Xi_{N_j}\Rightarrow\Xi$, the counts
$M_{N_j}(r)$ converge in law to $M(r)$ for each $r<1$, and
\begin{equation}\label{eq:moment-convergence}
 \E M_{N_j}(r)^q\longrightarrow\E M(r)^q,
 \qquad
 \E e^{tM_{N_j}(r)}\longrightarrow\E e^{tM(r)}.
\end{equation}
\end{theorem}

The proof has two separate parts. A radial comparison bounds the number
of particles in every compact window and gives tightness. A two-point
estimate then prevents different particles from merging in the limit.
Simplicity does not follow from tightness alone: for example,
$\delta_0+\delta_{1/n}\to2\delta_0$ vaguely, although each measure in
the sequence is simple.

\section{Radial comparison and compact-count tails}

We first prove the two elementary inequalities used in the comparison.

\begin{lemma}[Subharmonic means]\label{lem:subharmonic}
If $P$ is holomorphic and $\beta>0$, then $|P|^\beta$ is
subharmonic. Its circular mean about any center is nondecreasing in
the radius, for circles whose enclosed disks lie in its domain.
More generally, for a nonnegative continuous subharmonic function $u$,
\begin{equation}\label{eq:annular-mean}
 u(z)\le\frac{1}{\pi(b^2-a^2)}
 \int_{a<|w-z|<b}u(w)\,dA(w),\qquad 0\le a<b,
\end{equation}
provided the closed disk of radius $b$ about $z$ lies in the domain.
\end{lemma}
\begin{proof}
Put $\alpha=\beta/2$. Direct differentiation gives, for $\varepsilon>0$,
\[
 \Delta(|P|^2+\varepsilon)^\alpha
 =4\alpha(|P|^2+\varepsilon)^{\alpha-2}
       (\varepsilon+\alpha|P|^2)|P'|^2\ge0.
\]
The locally uniform limit as $\varepsilon\downarrow0$ is subharmonic.
For smooth $u$, the derivative of the circular mean at radius $s$ is
$(2\pi s)^{-1}\int_{B(z,s)}\Delta u\,dA\ge0$.
Regularization proves the corresponding statement for continuous
subharmonic functions. Integrate the circular submean inequality
against $2\pi s\dd s$ over $a<s<b$ to obtain
\eqref{eq:annular-mean}.
\end{proof}

\begin{lemma}[Increasing tilts of product measures]\label{lem:association}
Let $\mu$ be a product probability measure on finitely many intervals.
If $F\ge0$ is bounded, nondecreasing in each coordinate, and
$\mu(F)>0$, then the probability measure $\nu=F\mu/\mu(F)$ satisfies
\[
 \nu(A)\le\mu(A)
\]
for every coordinatewise decreasing Borel set $A$.
\end{lemma}
\begin{proof}
For bounded increasing functions $f,g$ of one variable and independent
identically distributed $X,X'$, one has
\[
 2\operatorname{Cov}(f(X),g(X))
 =\E[(f(X)-f(X'))(g(X)-g(X'))]\ge0.
\]
For a product measure, induct on the number of coordinates.
The conditional covariance given the last coordinate is nonnegative
by induction. The conditional expectations are increasing functions
of that last coordinate, so their covariance is also nonnegative.
The covariance decomposition proves the same inequality in every
dimension. Apply it to $f=F$ and $g=-\ind_A$ to obtain
$\mu(F\ind_A)\le\mu(F)\mu(A)$.
\end{proof}

By absolute continuity, zero or repeated radii have probability zero.
Sort the radii as $0<r_1<\cdots<r_N<1$ and relabel the angles
accordingly. Sorting contributes only a constant $N!$ to the density.
Define the ratios
\begin{equation}\label{eq:ratios}
 t_k=\frac{r_k}{r_{k+1}},\qquad r_{N+1}=1,
 \qquad r_j=\prod_{k=j}^N t_k,
 \qquad a_k=2k+\frac{\beta}{2}k(k-1).
\end{equation}
For $i<j$, factor
\[
 |z_i-z_j|=r_j\left|1-
   \left(\prod_{\ell=i}^{j-1}t_\ell\right)
        e^{i(\theta_i-\theta_j)}\right|.
\]
The triangular Jacobian and the radial Vandermonde factors give
\[
 \prod_{j=1}^N r_j\dd r_j=\prod_{k=1}^N t_k^{2k-1}\dd t_k,
 \qquad
 \prod_{j=1}^N r_j^{\beta(j-1)}
 =\prod_{k=1}^N t_k^{\beta k(k-1)/2}.
\]
Indeed, the Jacobian of $t\mapsto r$ is $\prod_k t_k^{k-1}$ and
$\prod_jr_j=\prod_kt_k^k$.
After angular integration, the ratio density is proportional to
\begin{equation}\label{eq:ratio-density}
 F_N(t)\prod_{k=1}^N t_k^{a_k-1},\qquad t\in(0,1)^N,
\end{equation}
where
\[
 F_N(t)=\int_{[0,2\pi)^N}
 \prod_{i<j}\left|1-
   \left(\prod_{\ell=i}^{j-1}t_\ell\right)
        e^{i(\theta_i-\theta_j)}\right|^\beta
       \prod_{j=1}^N\frac{\dd\theta_j}{2\pi}.
\]
This function is bounded, continuous, and nonnegative on $[0,1]^N$.

\begin{lemma}\label{lem:angular}
The angular factor $F_N$ is nondecreasing in every ratio $t_k$.
\end{lemma}
\begin{proof}
Fix $k<N$ and all other ratios. Only the pairs $i\le k<j$ depend
on $t_k$; their radial ratios have the form $t_kq_{ij}$.
Average additionally over the common rotation
$\theta_i\mapsto\theta_i+\phi$ for $i\le k$.
Invariance of angular Lebesgue measure leaves $F_N$ unchanged.
All noncrossing factors are unchanged, whereas the crossing product is
\[
 |P_\theta(t_ke^{i\phi})|^\beta,
 \qquad P_\theta(w)=\prod_{i\le k<j}
          (1-q_{ij}e^{i(\theta_i-\theta_j)}w).
\]
Its circular mean is nondecreasing in $t_k$ by
Lemma~\ref{lem:subharmonic}. Integration against the remaining
nonnegative factors preserves this inequality.
The factor $F_N$ does not depend on $t_N$.
\end{proof}

Let $U_k$, $1\le k\le N$, be independent with density
$a_ku^{a_k-1}$ on $(0,1)$. Equation~\eqref{eq:ratio-density} is
their product law tilted by $F_N$, whose mean is positive.
For $N\ge m$, the event $M_N(r)\ge m$ is the decreasing event
$\prod_{k=m}^N t_k<r$. Hence Lemma~\ref{lem:association} gives
\begin{equation}\label{eq:comparison}
 \Prob\{M_N(r)\ge m\}
 \le\Prob\left\{\sum_{k=m}^N E_k>\log\frac1r\right\},
 \qquad E_k=-\log U_k.
\end{equation}
The variables $E_k$ are independent exponentials with rates $a_k$:
$\Prob(E_k>x)=e^{-a_kx}$ for $x\ge0$.

For $m\ge2$, take $\lambda=\beta m(m-1)/4$.
For every $k\ge m$,
$a_k\ge\beta k(k-1)/2$ and $\lambda/a_k\le1/2$.
Using $-\log(1-x)\le2x$ for $0\le x\le1/2$,
\begin{align*}
 \log\E\exp\left(\lambda\sum_{k=m}^N E_k\right)
 &=\sum_{k=m}^N-\log(1-\lambda/a_k)\\
 &\le2\lambda\sum_{k=m}^N a_k^{-1}
 \le\frac{4\lambda}{\beta}
       \sum_{k=m}^\infty\frac1{k(k-1)}=m.
\end{align*}
Exponential Markov inequality and \eqref{eq:comparison} prove
\eqref{eq:tail} at finite $N$; if $N<m$, its left side is zero.
In particular, for every $r<1$, $q>0$, and $t>0$,
\begin{equation}\label{eq:uniform-moments}
 \sup_N\E M_N(r)^q<\infty,
 \qquad\sup_N\E e^{tM_N(r)}<\infty.
\end{equation}
For example, the terms with $m\ge2$ in the exponential-moment sum
are bounded by
$\exp((t+1)m-c_rm(m-1))$ with
$c_r=(\beta/4)\log(1/r)>0$, a summable sequence.

\section{Tightness in the space of random measures}

We recall the required topological facts and explain the deterministic
compactness and closure arguments. For a locally compact, second
countable space such as $\D$, $\Mplus(\D)$ with the vague topology
is Polish, meaning that its topology can be given by a complete
separable metric. We use Prokhorov's theorem: a family of probability
laws on a Polish space is tight if, for every $\varepsilon>0$, a
compact set has probability at least $1-\varepsilon$ under every law;
such a family is relatively sequentially compact for convergence in
law. For the Polish property and Prokhorov's theorem, see
\cite[Theorems A2.6.III and A2.4.I]{DV1}, respectively.

\begin{lemma}[Local mass bounds and counting-measure closure]
\label{lem:topology}
A family in $\Mplus(\D)$ whose masses are uniformly bounded on
each compact subset of $\D$ is relatively compact for the vague
topology. Moreover, $\Mp(\D)$ is closed in $\Mplus(\D)$.
\end{lemma}
\begin{proof}
For the first assertion, choose a compact exhaustion and countable
uniformly dense families of continuous test functions supported in
its successive members. The integrals of each test function are
bounded by the corresponding local mass bound. Diagonal extraction
gives convergence for this countable family. Uniform approximation
and the same mass bounds extend convergence to every function in
$C_c(\D)$. The limiting functional is positive and locally bounded,
so the Riesz representation theorem identifies it with a Radon
measure. Thus every sequence has a vaguely convergent subsequence.

For closure, suppose that counting measures $\xi_n$ converge vaguely
to $\mu\in\Mplus(\D)$. If a relatively compact ball $B$ satisfies
$\mu(\partial B)=0$, approximation of its indicator by continuous
functions gives $\xi_n(B)\to\mu(B)$. Therefore $\mu(B)$ is an
integer. For every $x\in\D$ there are arbitrarily small balls
centered at $x$ with $\mu$-null boundary: among radii below any fixed
compactly contained radius, only countably many circles can have
positive mass. Along decreasing such radii,
$\mu(B(x,s))\downarrow\mu(\{x\})$.
These ball masses are nonnegative integers and bounded above, so
they eventually stabilize. Thus $\mu(\{x\})$ is an integer, and
if it is zero some neighborhood of $x$ has zero mass.
It follows that every point of the support of $\mu$ is an atom of
positive integer mass. There are finitely many such atoms on each
compact set because $\mu$ is Radon. Hence $\mu$ is a locally finite
counting measure.
\end{proof}
The local compactness criterion and the closure assertion also appear
in \cite[Corollary A2.6.V]{DV1} and \cite[Lemma 9.1.V]{DV2},
respectively.

Choose $r_\ell\uparrow1$. Given $\varepsilon>0$, the tail bound
allows us to choose integers $b_\ell$ with
\[
 \sup_N\Prob\{\Xi_N(B_{r_\ell})>b_\ell\}
 <\varepsilon2^{-\ell}.
\]
Consider the subset of the full measure space
\[
 \mathcal K=\bigcap_{\ell\ge1}
 \{\mu\in\Mplus(\D):\mu(B_{r_\ell})\le b_\ell\}.
\]
Mass on an open relatively compact set is lower semicontinuous
for vague convergence. Thus $\mathcal K$ is closed.
Every compact subset of $\D$ lies in some $B_{r_\ell}$, so
Lemma~\ref{lem:topology} makes $\mathcal K$ relatively compact,
hence compact. A union bound gives
$\inf_N\Prob(\Xi_N\in\mathcal K)\ge1-\varepsilon$.
Prokhorov's theorem proves tightness in $\Mplus(\D)$.

Now fix any convergent subsequence $\Xi_{N_j}\Rightarrow\Xi$ in
that space. Since $\Mp(\D)$ is closed and contains every $\Xi_N$,
Portmanteau's theorem shows that $\Xi\in\Mp(\D)$ almost surely.
This excludes diffuse mass in the limit; multiplicities still require
proof.

The tail bound already passes to this counting-measure limit.
The set
$\{\mu:\mu(B_r)>m-1/2\}$ is open by lower semicontinuity. Thus
Portmanteau and the integer-valued counts give
\[
 \Prob\{\Xi(B_r)\ge m\}
 \le\liminf_j\Prob\{\Xi_{N_j}(B_r)\ge m\},
\]
which proves \eqref{eq:tail} for $\Xi$ and all its local moment
bounds. This argument requires no assumption about boundary points
or simplicity.

\section{A local two-point bound}

The finite-system one- and two-point densities, relative to $dA$
and $dA\otimes dA$, are defined by
\begin{align*}
 \E\sum_i f(z_i)&=\int_\D f(z)\rho_{1,N}(z)\,dA(z),\\
 \E\sum_{i\ne j}g(z_i,z_j)
 &=\int_{\D^2}g(z,w)\rho_{2,N}(z,w)\,dA(z)dA(w).
\end{align*}
The pairs in the second sum are ordered.
Integration of \eqref{eq:model} yields, for $p=1,2$ and $N\ge p$,
\begin{equation}\label{eq:correlation-factorization}
 \rho_{p,N}(z)=|\Delta_p(z)|^\beta h_{p,N}(z),
\end{equation}
where $(N)_p=N(N-1)\cdots(N-p+1)$ and
\begin{equation}\label{eq:h-definition}
 h_{p,N}(z)=\frac{(N)_p}{Z_{N,\beta}}
 \int_{\D^{N-p}}|\Delta_{N-p}(y)|^\beta
     \prod_{i=1}^p\prod_{\ell=1}^{N-p}|z_i-y_\ell|^\beta
           \prod_{\ell=1}^{N-p}dA(y_\ell).
\end{equation}
For $N<p$ set these functions to zero. For each fixed $N$, they
are continuous by dominated convergence. They are separately
subharmonic in their arguments, since the integrands are powers of
moduli of holomorphic polynomials, integrated against a positive
measure independent of $z$.

\begin{proposition}[Uniform local intensity and repulsion bounds]
\label{prop:repulsion}
For every $r<1$ there is $C_{\beta,r}<\infty$, independent of $N$,
such that, for $z,w\in\overline{B_r}$,
\begin{equation}\label{eq:repulsion}
 \rho_{1,N}(z)\le C_{\beta,r},\qquad
 \rho_{2,N}(z,w)\le C_{\beta,r}|z-w|^\beta.
\end{equation}
\end{proposition}
\begin{proof}
Fix $r<R<1$ and put $\eta=(R-r)/8$.
The disk submean inequality gives
\[
 h_{1,N}(z)\le\frac{1}{\pi\eta^2}
       \int_{B_R}\rho_{1,N}(u)\,dA(u)
 =\frac{\E M_N(R)}{\pi\eta^2}.
\]
This is uniformly bounded by \eqref{eq:uniform-moments}.

For $h_{2,N}$ we average its two arguments over separated regions.
If $|z-w|\ge4\eta$, use the disks $B(z,\eta)$ and $B(w,\eta)$;
points $u,v$ in these disks satisfy $|u-v|\ge2\eta$.
If $|z-w|<4\eta$, use $B(z,\eta)$ and the annulus
$6\eta<|v-w|<7\eta$. Then $|u-v|\ge\eta$.
Both choices lie inside $B_R$; each averaging region has area at
least $\pi\eta^2$. Successive applications of
Lemma~\ref{lem:subharmonic} therefore give in either case
\begin{align*}
 h_{2,N}(z,w)
 &\le\frac{\eta^{-\beta}}{\pi^2\eta^4}
       \int_{B_R^2}|u-v|^\beta h_{2,N}(u,v)\,dA(u)dA(v)\\
 &=\frac{\E[M_N(R)(M_N(R)-1)]}{\pi^2\eta^{\beta+4}}.
\end{align*}
The second moment bound is uniform in $N$.
Multiplication by $|z-w|^\beta$ proves the second assertion.
\end{proof}

\section{Why every accumulation point is simple}

Let $\Xi_{N_j}\Rightarrow\Xi$ be the subsequence fixed above.
We already know that $\Xi$ is a locally finite counting measure,
with multiplicities permitted. For such a measure $\xi$, its second
\emph{factorial measure} $\xi^{[2]}$ counts ordered pairs of distinct
particle labels, including pairs at the same location if an atom has
multiplicity. For $f\in C_c(\D^2)$,
\begin{equation}\label{eq:factorial-functional}
 \langle f,\xi^{[2]}\rangle
 =\int f(z,w)\,\dd\xi(z)\dd\xi(w)
       -\int f(z,z)\,\dd\xi(z).
\end{equation}
For example, an atom of mass $a$ contributes $a(a-1)$ ordered
pairs on the diagonal. This formula is valid before simplicity has
been established.

Both terms on the right of \eqref{eq:factorial-functional} are
continuous under vague convergence. For the product-measure term,
approximate $f$ uniformly on a compact product by finite sums of
products of one-variable continuous functions, with compactly
supported cutoffs. Vaguely convergent measures have uniformly bounded
mass on each compact set, so this approximation reduces the assertion
to the definition of vague convergence. The diagonal function
$z\mapsto f(z,z)$ is itself continuous with compact support.

If the coordinate projections of the support of $f$ lie in a compact
set $K\subset B_S$, then
\[
 |\langle f,\Xi_{N_j}^{[2]}\rangle|
 \le\|f\|_\infty M_{N_j}(S)^2.
\]
The uniform fourth moment in \eqref{eq:uniform-moments} makes these
random variables uniformly integrable. The continuous mapping theorem
and uniform integrability imply
\begin{equation}\label{eq:pair-moment-convergence}
 \E\langle f,\Xi_{N_j}^{[2]}\rangle
 \longrightarrow\E\langle f,\Xi^{[2]}\rangle.
\end{equation}
The same argument, using the second moment bound, gives
$\E\Xi_{N_j}(g)\to\E\Xi(g)$ for $g\in C_c(\D)$.

Write $\alpha_1=\E\Xi$ and $\alpha_2=\E\Xi^{[2]}$ for the
limiting intensity and second factorial moment measures.
They are locally finite by the count moment bounds.
For nonnegative continuous tests supported in $B_r$ or $B_r^2$,
Proposition~\ref{prop:repulsion} and
\eqref{eq:pair-moment-convergence} show that
\begin{equation}\label{eq:limiting-domination}
 \alpha_1\big|_{B_r}\le C_{\beta,r}\,dA,
 \qquad
 \alpha_2\big|_{B_r^2}
 \le C_{\beta,r}|z-w|^\beta\,dA(z)dA(w).
\end{equation}
These are inequalities of Radon measures: domination on nonnegative
continuous compactly supported tests extends to all Borel sets by
regularity. In particular, $\alpha_2$ gives no mass to the diagonal.
For every $r<1$,
\[
 0=\alpha_2\{(z,z):z\in B_r\}
   =\E\sum_{x\in B_r}
       \Xi(\{x\})\bigl(\Xi(\{x\})-1\bigr).
\]
The sum is nonnegative. It is therefore zero almost surely, so every
atom in $B_r$ has mass at most one. Taking a countable exhaustion
$r\uparrow1$ proves that $\Xi$ is simple almost surely.
This completes the assertion about accumulation points in the full
space of random measures.

The first inequality in \eqref{eq:limiting-domination} also proves
that $\Xi$ almost surely avoids every fixed area-null Borel subset
of $\D$. In particular, for every fixed $r<1$,
$\Xi(\partial B_r)=0$ almost surely. Counting $B_r$ is continuous
for vague convergence at any measure with this property: approximate
its indicator from inside and outside by continuous functions.
The continuous mapping theorem now gives
$M_{N_j}(r)\Rightarrow M(r)$.
Uniform bounds on higher powers and on $e^{sM_N(r)}$ for every $s>0$
give uniform integrability of each fixed power and exponential.
This proves \eqref{eq:moment-convergence} and finishes the proof of
Theorem~\ref{thm:main}.

\begin{remark}[Quantitative exclusion of close pairs]
The preceding proof gives a uniform estimate for both the finite
processes and every limit. If $K\subset B_r$ is compact and
$\varepsilon>0$, then
\[
 \Prob\{\text{two distinct points in }K
                    \text{ have distance less than }\varepsilon\}
 \le \frac{C_{\beta,r}\pi\,A(K)}{\beta+2}
               \varepsilon^{\beta+2},
\]
where $A(K)$ is the area of $K$. Indeed, such a pair contributes at
least two to the ordered pair count. Integrating the two-point bound
uses
$\int_{|u|<\varepsilon}|u|^\beta dA(u)
 =2\pi\varepsilon^{\beta+2}/(\beta+2)$.
\end{remark}

\begin{remark}[The role of positive inverse temperature]
At $\beta=0$ the points are independent and uniform on $\D$.
For $0<r<1$, $M_N(r)$ is binomial with parameters $N$ and $r^2$,
so it tends to infinity in probability. The local tightness result
therefore uses the assumption $\beta>0$ essentially.
\end{remark}

\begin{thebibliography}{2}
\bibitem{DV1}
D. J. Daley and D. Vere-Jones,
\emph{An Introduction to the Theory of Point Processes},
vol.~I: Elementary Theory and Methods, second edition, Springer, 2003.
Appendix A2.4--A2.6.

\bibitem{DV2}
D. J. Daley and D. Vere-Jones,
\emph{An Introduction to the Theory of Point Processes},
vol.~II: General Theory and Structure, second edition, Springer, 2008.
Section 9.1.
\end{thebibliography}
\end{document}
